\section{2 сентября 1914 г.}

\lp{20260604_193937725}

Швейцария. Цюрих. \\
Emile Medtner \\
Zurich \\
Hotel Bellevue

\lp{20260604_194010397}

Дорогой мой! Сегодня пришла ваша открытка от 27 Августа н.с. \rem{нового
стиля}, где вы пишите о том, чтобы я обратилась к Марго для «Мусагета».
Поговорю с ней. С Калей тоже поговорю. Всего я получила от вас 1 письмо, 2
телеграммы, 4 открытки. Коля и Каля по открытке, папаша письмо и Киселев. Я
пишу с 16 Августа старого стиля ежедневно, а отправляю через день, и все еще не
знаю получили ли вы хоть одно письмо.

Еще раз на всякий случай сообщаю вам, что Колю призывали, но освободили, т.к.
признали негодным. Все здоровы: у каждого из нас есть какая-нибудь забота,
имеющая отношение к событиям. Например, раненные. Сегодня мы с Верочкой делали
перевязку одному гренадеру.

У нас пока живет Ирочка. Коля ее готовил к экзамену. Теперь она начнет ходить в
консерваторию. Шура в Петербурге, а Боря пока еще не определился. Верочка
начала ходить в гимназию, но они там очень мало занимаются теперь науками.

Теперь буду посылать некоторое время открытки. До свидания, мой дорогой! Значит
на всю зиму?

Ваша Анюта.

\lp{20260604_194035089}

Швейцария. Цюрих. \\
Emile Medtner \\
Zurich \\
Hotel Bellevue

\lp{20260604_194117754}

Дорогой мой! Неужели не доходят до вас письма? Я послала их целую кучу. Начну
писать открытки. От вас получены: 2 телеграммы, 1 закрытое письмо, 3 открытки.
1 открытку получил Коля, Каля и письмо папаша. У нас все благополучно. Колю
признали негодным; я уже писала вам об этом. Живем, конечно, в напряжении.
Обнимаю! Ваша Анна.

\lp{20260604_193857242}

Швейцария. Цюрих. \\
Emile Medtner \\
Zurich \\
Hotel Bellevue

\lp{20260604_193920270}

Дорогой мой! 3 сентября к нам собирается Марго. Поговорю с ней о Мусогете.
Сегодня опустила вам открытку. У нас все благополучно. О деньгах не
беспокойтесь. Квартиру мы, конечно, не оставим. Писем я вам послала массу.
Обнимаю! Анюта
